% =============================
% ORGANIZAÇÃO DO CÓDIGO
% =============================

% ===
\section{Seção}
% ===

% ---
\subsection{Subseção}
% ---

%...
\subsubsection{Sub-subseção}
%...


% fim %

%%%%%%%%%%%%%%%%%%
Introduction
-----------------------------
Areas for Improvement:
-- Research Problem and Motivation: These subsections are placeholders. To strengthen the introduction, complete them with specific content. For example:
-- Research Problem: Highlight the persistence of environmental externalities in Brazil due to weak incentive structures, referencing specific policies (e.g., Código Florestal) and their limitations.
Motivation: Emphasize the urgency of the problem in Mato Grosso (e.g., deforestation rates, economic losses) and the inadequacy of current solutions (e.g., carbon markets’ limited adoption).
-- Hypotheses Overload: The 10 hypotheses are comprehensive but risk diluting focus. Consider consolidating related hypotheses (e.g., merge hypotheses 2 and 9 on carbon markets and biome-specific compensation) to streamline the research scope.
-- Objectives Alignment: The objectives are clear but could be more measurable. For example, “Establish a formal model” could specify “Develop a game-theoretic model of economic commensalism that predicts preservation rates under varying incentives.”
-- Include a Brazilian context earlier (e.g., mention Mato Grosso’s deforestation in the opening paragraph) to localize the global problem and justify the case study.
-- Refine the key question to be more specific, e.g., “How can market-based incentives in Brazil’s agricultural sector align private profitability with environmental preservation?”

-----------------------------
Hipóteses
-- Sugestão: Mantenha a Key Question como está ("How can economic incentives be structured to align private and social interests in environmental preservation?"). Para as Hypotheses, refine-as para serem mais específicas e testáveis, conectando-as diretamente ao economic commensalism. Por exemplo, a hipótese sobre carbono ("Efficiency of carbon markets") poderia ser reformulada como: "A structured carbon market, integrated into the economic commensalism framework, increases the adoption of sustainable practices by rural producers." Certifique-se de que cada hipótese seja testável nos capítulos de Mathematical Modeling e Empirical Validation.

- O ser humano é utilitário e não necessariamente bom.


%%%===========================================
\section{Literature Review}
Ver CP05052025a
-- Sugestão: A seção foi mantida como está, mas deve ser desenvolvida para ser analítica, não apenas descritiva. Aprofunde a discussão de Hardin, Pigou, Ostrom, e Georgescu-Roegen, identificando lacunas específicas (e.g., Pigou não lida com incerteza, Ostrom é limitada em escala, Georgescu-Roegen não oferece soluções práticas). Inclua outras perspectivas (e.g., Daly, Kapp, North) para reforçar a base teórica. Conecte cada lacuna ao economic commensalism, preparando o terreno para o Theoretical Framework.
-- Enfatizar os gaps da introduction e aprofundá-los, deixando claro o tamanho do problema

- Clássicos e neoclássicos

- Pigou, Hardin, Ostrom, Georgescu-Roegen: pilares históricos do pensamento ambiental econômico, com suas respectivas limitações sobre economia ambiental

Pigouvian taxes, Coasian bargaining, and Georgescu-Roegen’s sustainability principles.
Hardin’s classical framework and Ostrom’s institutional critique.

- North, Acemoglu: crítica institucional moderna — por que os incentivos ainda estão desalinhados.

- Assunção, Scheinkman: autores aplicados ao Brasil que demonstram que política ambiental depende de desenho inteligente e de monitoramento.

- Laffont, Daly, Martínez-Alier: reforços opcionais, usados pontualmente para apoiar partes específicas (como justiça distributiva ou falhas no desenho de políticas públicas).

Regarding externalities, this concept, introduced by \citeonline{pigou2010liberty}, added depth to this discourse by describing how economic agents' actions cause costs or benefits to others that are not compensated, such as pollution's impact on the environment. To rectify these market inefficiencies, Pigou suggested the implementation of corrective measures, like taxes or subsidies—referred to as Pigouvian taxes—to internalize the societal costs associated with externalities, thereby aligning personal incentives with societal well-being. Specifically, in the case of environmental harm, such taxes might deter excessive exploitation by ensuring that polluters account for the true costs of their activities.

However, subsequent contributions exposed limitations in the Pigovian approach. \citeonline{hardin1968tragedy} argued that the unregulated use of common-pool resources leads inevitably to overexploitation, given that each rational agent, when seeking to maximize their individual utility, ignores the collective repercussions of exploiting natural resources indiscriminately. For \citeonline{ostrom2002chapter24}, Hardin's assumption that economic agents are homogeneous in their utilitarian rationality may in fact occur, but primarily due to a lack of coordination and communication. Ostrom emphasizes that, with effective communication and institutional arrangements, appropriators can self-organize to manage resources sustainably, countering the tragedy of the commons.

Expanding on the critique of orthodox economics, Georgescu-Roegen challenged the traditional view of resource use as a closed, circular system, arguing that it ignores the entropic nature of economic processes. He posited that the exploitation of natural resources is inherently irreversible due to the degradation of energy and matter, generating unavoidable externalities such as pollution and resource depletion. For Georgescu-Roegen, sustainable resource management requires acknowledging these biophysical limits and adopting an evolutionary economic approach that prioritizes long-term ecological balance over short-term utility maximization, thereby aligning with Ostrom’s focus on adaptive, self-governed institutions to mitigate overexploitation \cite{cechinEconomiaEcologicaEvolucionaria2010}.

Fazer tabela
Conceito (comensalismo, extrativismo etc) | Paradigma teórico | solução atual (lei, incentivos etc)

PARETO, NASH E COMENSALISMO NA NATUREZA
o QUE A TESE FARÁ 

Legislação atual mira no mutualismo mas produz mais predatismo do que extrativismo?

\subsection{Game Theory and Equilibrium Analysis}
Nash equilibrium, correlated equilibrium, and strategic incentives for conservation.

\subsection{Behavioral Economics in Environmental Policy}
Bayesian learning, Markovian expectations, and the role of bounded rationality.
Não esquecer de coase
Gap conceitual

\newpage

% fim %

%%%===========================================
% ESTAVA EM SPECIFIC OBJECTIVES
\begin{enumerate}
    \item Establish a formal model incorporating \textit{economic commensalism} into production decisions.
    \item To develop a quantitative metric for evaluating the economic value of environmental resources, considering their contribution to sustainable production and carbon markets.
    \item Evaluate the effectiveness of positive (subsidies) and negative (fines, regulations) incentives in influencing conservation behaviors.
    \item Assess the feasibility of a localized versus globalized approach to carbon credit markets.
    \item Use econometric validation and historical case studies to test model predictions.
\end{enumerate}

- SOBRE RISCOS
a) Produtores rurais
* Maximizam lucro intertemporal.
* Têm informação imperfeita e enviesada sobre:
* a regeneração ambiental;
* o ponto de não retorno ecológico;
* a eficácia de políticas ambientais.
* São avessos a risco: o risco aqui é não obter lucro ou perder produtividade futura.

b) Governo/Estado
* Busca maximizar bem-estar social, com restrições orçamentárias e políticas.
* Teme os impactos negativos intertemporais e irreversíveis no meio ambiente.
* É averso ao risco político e ecológico: risco de colapso ambiental ou de conflito com produtores.

------------
- DEMONSTRAR CAP. 1 / ATACAR CAP. 2
a) Bellman
* Racionalidade intertemporal de ambos os agentes.

b) Kalman
* Ambos os agentes não observam diretamente o estado do meio ambiente.

b) Kalman
* Ambos os agentes não observam diretamente o estado do meio ambiente.
* Exemplo: área reflorestada ≠ área regenerada de fato.
* Aversão ao risco implica necessidade de estimar corretamente o “verdadeiro” estado latente.

c) Markov
* A evolução do meio ambiente (e da percepção sobre ele) é estocástica.
* A estrutura de decisão de ambos depende do estado corrente, não da história completa → propriedade de memória curta.

d) Spin
* As decisões de uso da terra (explorar ou preservar) são binárias, mas com efeitos interdependentes:
* Um produtor preservar pode afetar o incentivo do vizinho.
* Pode haver estados agregados de preservação (equilíbrios múltiplos).
* A aversão a risco e a pressão social podem gerar transições de estado em bloco (efeito manada ou transição de fase).

e) Bayesiano
* Como há incerteza sobre o meio ambiente e sobre a política pública futura, os agentes atualizam suas crenças:
* A aversão ao risco faz com que mesmo pequenos sinais negativos alterem significativamente o comportamento.

Hipóteses formais:
* A função utilidade dos agentes é côncava (representa aversão ao risco).
* O estado do meio ambiente é parcialmente observável.
* As decisões dos agentes dependem do estado corrente → Markov.
* O comportamento dos agentes é adaptativo → Bayes.
* As ações são binárias (explorar ou preservar) → Spin.

***
Chapter 1 sets the conceptual foundation for the thesis by addressing the motivation to endogenize the environment \( E \) within the production function. It begins with a critical review of the literature and the limitations of traditional economic models that treat nature as exogenous, that rely on externality correction mechanisms such as Pigouvian taxes, Coasian bargaining, and command-and-control approaches. Building on this, it argues for incorporating ecological concepts—predatism, mutualism, extractivism, and especially \textit{commensalism}—as heuristic tools to reinterpret economic-environmental interactions. The chapter then explores how introducing \( E \) into isoquant and isocost analysis reshapes marginal trade-offs and enables a new theoretical framework. This leads to the formal definition of \textit{economic commensalism} as a neutral-impact paradigm, offering an alternative to both exploitative and overly restrictive models of environmental use.

% ==============================
\textbf{Preliminary title}: 

This chapter formalizes the hypotheses presented in Chapter 1 through a dynamic optimization model that integrates Bellman recursion, Markov uncertainty, Kalman filtering, Bayesian updating, and binary decision dynamics inspired by Spin networks. The model aims to demonstrate under what conditions an agent — facing uncertainty and learning over time — may rationally choose preservation over exploitation.

\textbf{Methodology}: 

\textbf{Justification}: This model 
